\section{Discussion}
\label{sec:discussion}

If the handover is anticipated, the push phase moves most of the work to the time before the stop. For pre-copy this is the transfer, for manifest-push the install. As derived (RQ-1), the push phase shortens the interruption time in the live study. Part of this reduction comes from creating the target before the stop. However, pre-copy has a longer total handover time than the full copy. Both correspond to Clark et al., who reserve the resources on the target before the copy and describe how to balance the downtime against the total migration time \cite{clark2005live}. Furthermore, the push phase of pre-copy puts extra work on the source, which shows in its source shortfall (Section~\ref{sec:results}). For a source close to its limit, this speaks in favor of manifest-push.

A single push round is enough in this setup. In all 60 runs with a push phase, only three small files follow at the stop. Two of them, the log and the state file, form the small writable working set of the heartbeat (Section~\ref{sec:related-work}), which makes it an easy case for pre-copy \cite{clark2005live}. Further rounds wouldn't help with the heartbeat, because it keeps writing to its files, so the set of changed files doesn't shrink.

Without anticipation, the choice between the full copy and manifest-based transfer depends on the project size (RQ-2). At the large level, manifest-based transfer is faster mainly because it moves most of the transfer from the uplink of the client to the registry link of the target. Therefore, this advantage depends on the registry link, which isn't varied in this setup.

At the small level, both strategies still take about \SI{3}{\second} in the median (Table~\ref{tab:interruption}). This size-independent share of E2B in us-west-1 applies to strategies that create the target after the stop, so the interruption doesn't grow proportionally with the project size.

Part of this share is the network path to the US. A target closer to the client would shorten it. The shortest round trip to the target has a median of 192 to \SI{195}{\milli\second} in every cell. Within the interruption, the runner makes four calls to E2B one after another for the full copy, five for manifest-based transfer with an install and three for the push strategies. For comparison, gcping measured an HTTP latency of \SI{195}{\milli\second} from Northern Germany to Google Cloud us-west1, where the default region of E2B runs \cite{e2b-egress-2026}, and \SI{41}{\milli\second} to europe-west3 \cite{gcping-2026}.

\subsection{Threats to Validity}

The heartbeat neither loads nor writes its dependencies. Therefore, a successful install only means that the install command ends without an error. Furthermore, the dependencies never change on the source during the push phase, which both push strategies assume. The heartbeat counter is the only process state that crosses the boundary, because the heartbeat writes it to a file. Since there is no memory image, state such as open connections or the module graph of a development server would be lost at every handover. Lastly, the full copy and pre-copy send the dependencies installed in WebContainers unchanged to the target, and none of the three packages carries binaries for a specific platform.

Only one client is used, and each study runs in one measurement session with one uplink, one region and one hosted target. Furthermore, the load on the E2B hosts isn't under control. For manifest-based transfer, the bytes transferred only count the archive and not what the install downloads from the registry. The handover is triggered by the script, as the WebContainers \ac{API} documents no signal before a resource runs out \cite{stackblitz-2024}. Hence, detecting that a resource runs out isn't part of the interruption time.

Within this setup, the crossing point depends on four further factors. First, each upper level carries a different direct dependency. Hence, the effect of the size can't be separated from the effect of the packages. Second, every megabyte adds about \SI{242.8}{\milli\second} of upload time over the 120 runs of the full copy. If the uplink were half as fast, the crossing point would move to \SI{5.9}{\mega\byte}, and if it were twice as fast, to \SI{13.7}{\mega\byte}. Third, the phases on the client lie inside the interruption time. For the full copy at the large level, they take 16.3\,\% of the median interruption time, almost all of it for the export. For manifest-based transfer, they take less than 1\,\%. Therefore, a weaker client would move the crossing point towards smaller projects. Lastly, the straight line between the medium and the large level was chosen after the measurement session. Other estimates from the same data compare the transfer, the install and the target start instead of the interruption time. In these estimates, the transfer of the full copy grows with the payload on a straight line over its 120 runs. The install of manifest-based transfer is taken from one of the three upper levels. Depending on this level, the crossing point lies between 5.8 and \SI{10.3}{\mega\byte}.

A Bonferroni correction across the four levels raises the bound of $p$ to 0.0004 and changes none of the results. The interruption time also depends on reading the clock of the target (Section~\ref{sec:methodology}). The bound of the reading error, half of the shortest round trip of each run, is at most \SI{337}{\milli\second}. Therefore, the difference between two strategies changes by at most \SI{674}{\milli\second}, which is smaller than the differences at the three upper levels and in the primary comparisons.
